%!TEX root=../mythesis.tex
% Appendix A

\chapter{Proofs and Derivations} % Main appendix title
\label{AppendixA}

This appendix collects the formal results referenced from the technical chapters: the convergence of the EarnHFT Q-teacher (\sref{app:earnhft_proof}), the convergence of the FineFT selective update (\sref{app:fineft_proof}), and the generalisation properties of the AlphaQD formulation (\sref{app:alphaqd_theory}).

%----------------------------------------------------------------------------------------
\section{Convergence of the EarnHFT Q-Teacher}
\label{app:earnhft_proof}

We prove Theorem~\ref{thm:earnhft_convergence} in its tabular idealisation: on the historical-replay MDP whose state is the pair $(t,p)$ of timestamp and position---the space on which the teacher table of Algorithm~\ref{alg:earnhft_qtable} is defined---with discount $\gamma\in(0,1)$, the action value learned with the idealised supervised update converges almost surely to $Q^\star$, provided the step sizes satisfy the Robbins--Monro conditions ($\sum_k a_k=\infty$, $\sum_k a_k^2<\infty$), every state--action pair is visited infinitely often, and the update noise has bounded variance. The analysis concerns the linear advantage-matching operator below, not the KL-form penalty of \eqref{eq:earnhft_loss} itself; no formal correspondence between the decaying coefficient $\alpha_{\mathrm{KL}}$ of \eqref{eq:earnhft_loss} and the fixed mixing weight $\lambda$ below is established. The method is based on $Q$-learning, so following the standard convergence argument~\cite{jaakkola1993convergence} it suffices, under the stated stochastic-approximation conditions, to show that the supervised update operator $\mathcal{T}$ is a contraction in the sup-norm whose fixed point is $Q^\star$. The operator combines the usual Bellman backup with the supervision term and is
\begin{equation}
\label{eq:app_earnhft_operator}
\mathcal{T}q(x,a) = \lambda\sum_{y\in X}T(y\mid x,a)\bigl(r(x,a,y)+\gamma\max_b q(y,b)\bigr) + (1-\lambda)\Bigl(\tfrac{1}{|A|}\sum_{a\in A}q(x,a)+Re(x,a)\Bigr),
\end{equation}
where the supervision residual is
\begin{equation}
\label{eq:app_earnhft_residual}
Re(x,a) = q^\star(x,a) - \tfrac{1}{|A|}\sum_{a\in A}q^\star(x,a),
\end{equation}
and $\lambda\in(0,1]$ weights the Bellman backup against the supervision term.

\begin{proposition}[$Q^\star$ is a fixed point]
\label{prop:app_earnhft_fixed}
The optimal action value $q^\star$ satisfies $\mathcal{T}q^\star = q^\star$.
\end{proposition}
\begin{proof}
Substituting $q^\star$ into \eqref{eq:app_earnhft_operator} and using the Bellman optimality equation $q^\star(x,a)=\sum_y T(y\mid x,a)(r(x,a,y)+\gamma\max_b q^\star(y,b))$ for the first term and \eqref{eq:app_earnhft_residual} for the second,
\begin{align*}
\mathcal{T}q^\star(x,a) &= \lambda\, q^\star(x,a) + (1-\lambda)\Bigl(\tfrac{1}{|A|}\textstyle\sum_a q^\star(x,a) + q^\star(x,a) - \tfrac{1}{|A|}\textstyle\sum_a q^\star(x,a)\Bigr) \\
&= \lambda\, q^\star(x,a) + (1-\lambda)\, q^\star(x,a) = q^\star(x,a). \qedhere
\end{align*}
\end{proof}

\begin{proposition}[Contraction in the sup-norm]
\label{prop:app_earnhft_contraction}
Let $\gamma\in(0,1)$. For any action values $q_1,q_2$, $\|\mathcal{T}q_1-\mathcal{T}q_2\|_\infty \le k\,\|q_1-q_2\|_\infty$ with $k=\lambda\gamma+(1-\lambda)=1-\lambda(1-\gamma)<1$.
\end{proposition}
\begin{proof}
The supervision residual $Re(x,a)$ depends only on $q^\star$, not on the argument of $\mathcal{T}$, so it cancels in the difference $\mathcal{T}q_1-\mathcal{T}q_2$. Writing the difference and bounding the $\max$ operator by $|\max_b q_1(y,b)-\max_b q_2(y,b)|\le\max_b|q_1(y,b)-q_2(y,b)|$,
\begin{align*}
\|\mathcal{T}q_1-\mathcal{T}q_2\|_\infty
&= \Bigl\| \lambda\textstyle\sum_y T(y\mid x,a)\gamma\bigl(\max_{b}q_1(y,b)-\max_{b}q_2(y,b)\bigr) \\
&\qquad + (1-\lambda)\tfrac{1}{|A|}\textstyle\sum_{a'}\bigl(q_1(x,a')-q_2(x,a')\bigr)\Bigr\|_\infty \\
&\le \lambda\gamma\,\|q_1-q_2\|_\infty + (1-\lambda)\,\|q_1-q_2\|_\infty \\
&= \bigl(\lambda\gamma+(1-\lambda)\bigr)\|q_1-q_2\|_\infty = \bigl(1-\lambda(1-\gamma)\bigr)\|q_1-q_2\|_\infty.
\end{align*}
Since $\gamma\in(0,1)$ and $\lambda\in(0,1]$, the factor $k=1-\lambda(1-\gamma)<1$, so $\mathcal{T}$ is a contraction. When $\gamma=1$ the factor equals $1$ and $\mathcal{T}$ is merely non-expansive; the undiscounted case is instead handled directly by the finite-horizon backward induction of Algorithm~\ref{alg:earnhft_qtable}.
\end{proof}

Combining the two propositions, $\mathcal{T}$ is a contraction with unique fixed point $q^\star$; under the Robbins--Monro step-size, infinite-visitation, and bounded-noise-variance conditions stated above, the stochastic-approximation argument of~\cite{jaakkola1993convergence} applies and the iterates converge almost surely to $Q^\star$, which proves Theorem~\ref{thm:earnhft_convergence}. The deployed implementation, which uses function approximation, lies outside this tabular idealisation; the result should be read as motivation rather than a guarantee for the deployed algorithm.

%----------------------------------------------------------------------------------------
\section{Convergence of the FineFT Selective Update}
\label{app:fineft_proof}

We argue that the selective-update rule of \sref{sec:fineft_stage1} preserves convergence of each learner on the dynamic it specialises in. The argument has two parts.

First, the pre-training loss \eqref{eq:fineft_pretrain} updates every learner with the same optimal-value supervisor used by EarnHFT (now constrained by the liquidation mask of Algorithm~\ref{alg:fineft_coav}). The per-learner update operator is exactly the operator $\mathcal{T}$ of \eqref{eq:app_earnhft_operator}, so by \sref{app:earnhft_proof} each learner's action value converges, in the same tabular idealisation and under the same stochastic-approximation conditions, to the constrained optimal action value on the demonstration distribution.

Second, consider the selective-update phase restricted to transitions from a single dynamic $z$. By the DP-MDP assumption (\sref{sec:fineft_dpmdp}), these transitions are generated by a stationary MDP $M_z$. The weight matrix of Algorithm~\ref{alg:fineft_weight} assigns, for each transition, a non-negative weight to the best-fitting learner $i^\star$ and its neighbours, and zero elsewhere; the diagonal weight of $i^\star$ is one. The update of learner $i^\star$ on transitions of dynamic $z$ is therefore a standard Huber-TD update under $M_z$ with a non-negative supervision term, whose update operator is a sup-norm contraction by the same argument as \sref{app:earnhft_proof}. The positive-feedback property---that specialising in $z$ lowers learner $i^\star$'s TD error on $z$, which raises its selection weight on subsequent $z$-transitions---ensures that, once a learner becomes the best fit for $z$, it continues to receive the $z$-transitions that drive its convergence on $M_z$. The cross-learner weights of neighbouring learners are bounded by the diagonal weights and decay with index distance, so they perturb the update by a contraction-preserving convex combination and do not break the contraction. Hence each specialist converges to the optimal action value of the dynamic it is selected for.

The full formal treatment, including the precise stochastic-approximation conditions on the selection weights, follows the convergence analysis of the original FineFT manuscript.

%----------------------------------------------------------------------------------------
\section{Generalisation Properties of AlphaQD}
\label{app:alphaqd_theory}

This section gives the full statements and proofs of the results summarised in \sref{sec:alphaqd_theory}. Two assumptions are used by name. \textbf{(A1)} \emph{Sign-symmetric warm-up shape distribution}: the warm-up sample $\{\mathbf{r}(\alpha_i)\}_{i=1}^M$ has zero mean and a second moment equal to that of its sign-augmented version $\{\mathbf{r}(\alpha_i),-\mathbf{r}(\alpha_i)\}$, which holds because a formula and its sign-flip are both legal under the same mask and produce IC series differing only by a global sign. \textbf{(A2)} \emph{Positive ridge weights on retained orbit pairs}: the shared ridge combiner assigns strictly positive weight to both representatives of an orbit when both are exported, which holds in the empirical regime where $\ell_2$-regularised regression on positively-edged factors places weight on both inputs.

\subsection{Decoupling Quality from Shape}
\label{app:alphaqd_decouple}

\begin{lemma}[Affine invariance of the shape descriptor; restates Lemma~\ref{lem:alphaqd_affine}]
\label{app:lem_affine}
For any candidate $\alpha$ and any positive affine transform $\mathbf{y}=a\,\mathbf{x}(\alpha)+b\mathbf{1}$ with $a>0$, the normalised shape vector built from $\mathbf{y}$ equals $\mathbf{r}(\alpha)$.
\end{lemma}
\begin{proof}
The mean satisfies $\overline{\mathbf{y}}=a\,\overline{\mathbf{x}(\alpha)}+b$, so $\mathbf{y}-\overline{\mathbf{y}}\mathbf{1}=a(\mathbf{x}(\alpha)-\overline{\mathbf{x}(\alpha)}\mathbf{1})$, and normalising by the $\ell_2$ norm cancels $a>0$, giving $\mathbf{r}(\alpha)$.
\end{proof}

The consequence is that the descriptor encodes only \emph{which} dates were relatively strong, not the overall edge: mean and scale are absorbed into the ICIR, removing the leakage that would let two alphas with the same temporal pattern but different edge occupy different cells.

\subsection{Reward Stationarity}
\label{app:alphaqd_stationary}

\begin{theorem}[Stationary conditional reward; restates Theorem~\ref{thm:alphaqd_stationary}]
\label{app:thm_stationary}
Let $R_q(\alpha)=\exp(\beta_q\,\mathrm{ICIR}(\alpha))$. For any fixed target orbit $\mathcal{O}^\star$ and any kernel $K:\mathcal{O}\times\mathcal{O}\to\mathbb{R}_{\ge 0}$ independent of the archive state and the training-step index, the conditional reward $R_{\mathcal{O}^\star}(\alpha)=R_q(\alpha)\,K(\mathcal{O}(c(\alpha)),\mathcal{O}^\star)$ is a stationary terminal reward.
\end{theorem}
\begin{proof}
For each fixed $\mathcal{O}^\star$, $R_{\mathcal{O}^\star}(\cdot)$ depends on $\alpha$ only through $\mathrm{ICIR}(\alpha)$ on the fixed training window and through the fixed PCA basis and bin edges entering $c(\cdot)$. The kernel $K$ is fixed and $\mathcal{O}^\star$ is a parameter, not a function of training-time state, so $R_{\mathcal{O}^\star}$ has no time-dependent component. A scheduler that varies $\mathcal{O}^\star$ across episodes changes only the per-episode prior over conditions, never the conditional reward.
\end{proof}

The implementation uses the unconditional reward ($K\equiv 1$), so the property holds a fortiori; the conditional form documents compatibility with future descriptor-biased exploration extensions.

\subsection{Search Complexity and the Boosting Counter-Example}
\label{app:alphaqd_complexity}

\begin{theorem}[Residual-edge noise floor]
\label{app:thm_boost}
Let $\mathcal{F}=\{\alpha_1,\dots,\alpha_M\}$ be a finite candidate set. After $t-1$ boosting rounds with residual $r_{t-1}$, the empirical residual-edge maximiser is $\alpha_t=\argmax_{\alpha\in\mathcal{F}}\widehat{\langle\alpha,r_{t-1}\rangle}$. Let $\gamma_t(\alpha)=\mathbb{E}[\langle\alpha,r_{t-1}\rangle]$ be the population residual edge, $\Delta_t=\sup_\alpha\gamma_t(\alpha)$, and $\eta_M=\sigma\sqrt{2\log(2M/\delta)/n_{\mathrm{eff}}}$ a sub-Gaussian deviation. Then with probability at least $1-\delta$, $\gamma_t(\alpha_t)\ge\Delta_t-2\eta_M$; when $\Delta_t\le 2\eta_M$, the bound is non-positive.
\end{theorem}
\begin{proof}
Uniform sub-Gaussian concentration over $\mathcal{F}$ gives $\sup_\alpha|\hat\gamma_t(\alpha)-\gamma_t(\alpha)|\le\eta_M$ with probability $\ge 1-\delta$. Writing $\alpha_t^\star=\argmax_\alpha\gamma_t(\alpha)$, $\gamma_t(\alpha_t)\ge\hat\gamma_t(\alpha_t)-\eta_M\ge\hat\gamma_t(\alpha_t^\star)-\eta_M\ge\gamma_t(\alpha_t^\star)-2\eta_M=\Delta_t-2\eta_M$, where the middle step uses the empirical optimality of $\alpha_t$.
\end{proof}

With $M>10^{20}$ and $n_{\mathrm{eff}}$ bounded by training-window days, the late-round residual winner is governed by the noise floor $\eta_M$ rather than by population signal.

\begin{theorem}[Search-complexity contrast; restates Theorem~\ref{thm:alphaqd_complexity}]
\label{app:thm_complexity}
A $K$-step boosting procedure over $\mathcal{F}$ with $|\mathcal{F}|=M$ has an ordered-sequence hypothesis class of size at most $M^K$, with finite-class complexity $\mathcal{C}_{\mathrm{boost},K}=\sqrt{K\log M/n_{\mathrm{eff}}}$. A quality-diversity rule mapping candidates to at most $G$ occupied cells and exporting one representative per cell has a $K$-subset class of size at most $\binom{G}{K}$, with complexity $\mathcal{C}_{\mathrm{QD},K}=\sqrt{K\log(G/K)/n_{\mathrm{eff}}}$. For $G\ll M$, $\mathcal{C}_{\mathrm{QD},K}\ll\mathcal{C}_{\mathrm{boost},K}$.
\end{theorem}
\begin{proof}
The finite-class generalisation bound scales as $\sqrt{\log|\mathcal{H}|/n_{\mathrm{eff}}}$. Substituting $|\mathcal{H}_{\mathrm{boost},K}|\le M^K$ gives $\log|\mathcal{H}_{\mathrm{boost},K}|\le K\log M$. Substituting $|\mathcal{H}_{\mathrm{QD},K}|\le\binom{G}{K}$ and $\log\binom{G}{K}\le K\log(eG/K)$ gives the QD scaling.
\end{proof}

In the reported grid $G\le 100$ while $|\mathcal{A}|>10^{20}$, so the gap spans several orders of magnitude.

\subsection{Per-Cell Selection Bias}
\label{app:alphaqd_selection}

\begin{theorem}[Per-cell selection bias under a bounded quota]
\label{app:thm_selection}
Fix an occupied cell holding at most $\kappa$ admitted elites at selection time, each with population quality $\theta_i$ and training-window estimate $\hat\theta_i=\theta_i+\varepsilon_i$ with sub-Gaussian $\varepsilon_i$ of effective sample size $n_{\mathrm{eff}}$. With probability at least $1-\delta$, the selection bias of the per-cell representative on raw ICIR satisfies $\max_{i\le\kappa}(\hat\theta_i-\theta_i)\le\sigma\sqrt{2\log(\kappa/\delta)/n_{\mathrm{eff}}}$. If the ranking uses the PFS-augmented score averaging the ICIR with $m$ independent perturbations, $\sigma$ is replaced by $\sigma/\sqrt{1+m}$.
\end{theorem}
\begin{proof}
A union bound on the $\kappa$ sub-Gaussian tails gives $\mathbb{P}(\max_{i\le\kappa}\varepsilon_i>t)\le\kappa\exp(-n_{\mathrm{eff}}t^2/(2\sigma^2))$; setting the right-hand side to $\delta$ yields $t=\sigma\sqrt{2\log(\kappa/\delta)/n_{\mathrm{eff}}}$. Averaging $1+m$ independent sub-Gaussian estimates with parameter $\sigma$ yields parameter $\sigma/\sqrt{1+m}$.
\end{proof}

The bounded quota replaces the candidate-count exponent that would govern an unbounded per-cell maximum by the constant $\kappa$; under $\kappa=3$, $\delta=0.05$ the bound is $\approx 2.86\,\sigma/\sqrt{n_{\mathrm{eff}}}$, and the $m=5$ perturbations tighten it by $\sqrt{6}$ when the PFS-augmented score decides the ranking.

\subsection{Ensemble ICIR Lower Bound}
\label{app:alphaqd_ensemble}

\begin{theorem}[Ensemble ICIR lower bound; restates Theorem~\ref{thm:alphaqd_ensemble}]
\label{app:thm_ensemble}
If the export retains $K$ positive-ICIR elites with daily IC variables $X_1,\dots,X_K$ satisfying $\mathbb{E}[X_i]=\mu_i\ge\mu_{\min}>0$, $\Var(X_i)\le\sigma^2$, and $\corr(X_i,X_j)\le\rho_{\max}$, then the equal-weight ensemble $X_{\mathrm{ens}}=\frac{1}{K}\sum_i X_i$ obeys $\mathrm{ICIR}(X_{\mathrm{ens}})\ge\mu_{\min}\sqrt{K}/(\sigma\sqrt{1+(K-1)\rho_{\max}})$.
\end{theorem}
\begin{proof}
By linearity, $\mathbb{E}[X_{\mathrm{ens}}]\ge\mu_{\min}$. For the variance,
\[
\Var(X_{\mathrm{ens}}) = \tfrac{1}{K^2}\textstyle\sum_i\Var(X_i) + \tfrac{1}{K^2}\sum_{i\ne j}\Cov(X_i,X_j) \le \tfrac{\sigma^2}{K^2}\bigl[K+K(K-1)\rho_{\max}\bigr] = \tfrac{\sigma^2(1+(K-1)\rho_{\max})}{K},
\]
so $\mathrm{ICIR}(X_{\mathrm{ens}})=\mathbb{E}[X_{\mathrm{ens}}]/\sqrt{\Var(X_{\mathrm{ens}})}\ge\mu_{\min}\sqrt{K}/(\sigma\sqrt{1+(K-1)\rho_{\max}})$.
\end{proof}

The bound rests on positive standalone edges, enforced by the positive-ICIR admission rule, and on controlled redundancy $\rho_{\max}$, enforced by the bounded per-cell quota and the PCA descriptor; \sref{sec:alphaqd_diagnostics} reports an empirical pairwise correlation of $0.2478$. Theorem~\ref{app:thm_boost} shows that boosting-style fitting preserves neither premise, so no analogous bound is available for it.

%----------------------------------------------------------------------------------------
\section{Full AlphaQD Ablation Grid}
\label{app:alphaqd_ablation}

\tref{tab:app_ablation_full} reports the full $36$-configuration ablation grid summarised by the marginal means in \sref{sec:alphaqd_ablation}. All entries are evaluated on the A-share test window under the unified-ridge protocol with a single training seed and the canonical training budget ($2{,}000$ episodes). The only varied hyperparameters are descriptor dimension $d_b$, cell capacity $\kappa$, and reward temperature $\beta_q$. Within-grid noise (ICIR roughly $0.22$--$0.66$) is larger than the marginal-mean effects reported in \sref{sec:alphaqd_ablation}, so individual grid cells should be read as samples from the AlphaQD performance distribution rather than as separated comparisons.

\begin{table}[!htb]
\centering
\caption[Full AlphaQD ablation grid.]{AlphaQD ablation: per-configuration metrics on the A-share panel under the unified-ridge protocol (single seed, test window). Columns are descriptor dimension $d_b$, cell capacity $\kappa$, reward temperature $\beta_q$, signal IC, ICIR, Spearman rank-IC, and net Sharpe.}
\label{tab:app_ablation_full}
\renewcommand{\arraystretch}{1.05}
\footnotesize
\begin{tabular}{ccccccc}
\toprule
$d_b$ & $\kappa$ & $\beta_q$ & IC & ICIR & RankIC & Sharpe \\
\midrule
$1$ & $1$ & $0.5$ & $0.0565$ & $0.455$ & $0.074$ & $0.322$ \\
$1$ & $1$ & $1.0$ & $0.0584$ & $0.592$ & $0.097$ & $0.951$ \\
$1$ & $1$ & $2.0$ & $0.0772$ & $0.647$ & $0.127$ & $1.027$ \\
$1$ & $1$ & $5.0$ & $0.0328$ & $0.246$ & $0.045$ & $0.889$ \\
$1$ & $3$ & $0.5$ & $0.0740$ & $0.542$ & $0.123$ & $1.186$ \\
$1$ & $3$ & $1.0$ & $0.0577$ & $0.464$ & $0.094$ & $0.699$ \\
$1$ & $3$ & $2.0$ & $0.0757$ & $0.630$ & $0.106$ & $0.540$ \\
$1$ & $3$ & $5.0$ & $0.0760$ & $0.511$ & $0.127$ & $0.835$ \\
$1$ & $5$ & $0.5$ & $0.0669$ & $0.493$ & $0.096$ & $0.833$ \\
$1$ & $5$ & $1.0$ & $0.0658$ & $0.560$ & $0.100$ & $0.879$ \\
$1$ & $5$ & $2.0$ & $0.0683$ & $0.664$ & $0.107$ & $1.010$ \\
$1$ & $5$ & $5.0$ & $0.0486$ & $0.485$ & $0.077$ & $0.668$ \\
\midrule
$2$ & $1$ & $0.5$ & $0.0541$ & $0.461$ & $0.075$ & $0.874$ \\
$2$ & $1$ & $1.0$ & $0.0258$ & $0.218$ & $0.040$ & $0.909$ \\
$2$ & $1$ & $2.0$ & $0.0724$ & $0.662$ & $0.106$ & $1.175$ \\
$2$ & $1$ & $5.0$ & $0.0760$ & $0.646$ & $0.096$ & $0.738$ \\
$2$ & $3$ & $0.5$ & $0.0736$ & $0.537$ & $0.123$ & $0.955$ \\
$2$ & $3$ & $1.0$ & $0.0432$ & $0.366$ & $0.047$ & $0.907$ \\
$2$ & $3$ & $2.0$ & $0.0724$ & $0.541$ & $0.112$ & $0.718$ \\
$2$ & $3$ & $5.0$ & $0.0712$ & $0.565$ & $0.113$ & $0.869$ \\
$2$ & $5$ & $0.5$ & $0.0589$ & $0.529$ & $0.087$ & $0.976$ \\
$2$ & $5$ & $1.0$ & $0.0793$ & $0.613$ & $0.122$ & $0.892$ \\
$2$ & $5$ & $2.0$ & $0.0375$ & $0.511$ & $0.087$ & $1.005$ \\
$2$ & $5$ & $5.0$ & $0.0734$ & $0.578$ & $0.123$ & $1.141$ \\
\midrule
$3$ & $1$ & $0.5$ & $0.0699$ & $0.505$ & $0.109$ & $0.991$ \\
$3$ & $1$ & $1.0$ & $0.0318$ & $0.368$ & $0.058$ & $0.669$ \\
$3$ & $1$ & $2.0$ & $0.0502$ & $0.516$ & $0.080$ & $0.644$ \\
$3$ & $1$ & $5.0$ & $0.0669$ & $0.534$ & $0.083$ & $1.144$ \\
$3$ & $3$ & $0.5$ & $0.0756$ & $0.525$ & $0.127$ & $1.040$ \\
$3$ & $3$ & $1.0$ & $0.0603$ & $0.518$ & $0.078$ & $0.581$ \\
$3$ & $3$ & $2.0$ & $0.0554$ & $0.616$ & $0.085$ & $0.957$ \\
$3$ & $3$ & $5.0$ & $0.0355$ & $0.405$ & $0.048$ & $-0.008$ \\
$3$ & $5$ & $0.5$ & $0.0314$ & $0.352$ & $0.031$ & $0.422$ \\
$3$ & $5$ & $1.0$ & $0.0741$ & $0.562$ & $0.130$ & $1.025$ \\
$3$ & $5$ & $2.0$ & $0.0494$ & $0.414$ & $0.087$ & $0.584$ \\
$3$ & $5$ & $5.0$ & $0.0306$ & $0.342$ & $0.044$ & $0.847$ \\
\bottomrule
\end{tabular}
\end{table}
